\section{Exact reconstruction from structural zeros}
\label{sec:fourier-reconstruction}

Let \(G=\Z/q\Z\), where \(q\ge2\), and write
\[
 e_q(x)=\ee^{2\pi i x/q}.
\]
For \(A:G\to\C\), use the unnormalized Fourier transform
\begin{equation}
 \widehat A(r)=\sum_{n\in G}A(n)e_q(-rn),
 \qquad r\in G.
\label{eq:dft}
\end{equation}
Then
\begin{equation}
 A(n)=\frac1q\sum_{r\in G}\widehat A(r)e_q(rn),
\qquad
 \sum_{r\in G}|\widehat A(r)|^2
 =q\sum_{n\in G}|A(n)|^2.
\label{eq:inversion-parseval}
\end{equation}
Put
\[
 Z=\widehat A(0)=\sum_nA(n),\qquad
 D=\sum_n|A(n)|^2.
\]

\begin{proposition}[One-zero coherent reconstruction]
\label{prop:one-zero}
If \(A(k)=0\) for some \(k\in G\), then
\begin{equation}
 \boxed{
 Z=-\sum_{r\ne0}\widehat A(r)e_q(rk).}
\label{eq:one-zero}
\end{equation}
In particular, if \(A(0)=0\), then
\begin{equation}
 Z=-\sum_{r\ne0}\widehat A(r).
\label{eq:zero-coordinate}
\end{equation}
\end{proposition}

\begin{proof}
Insert \(n=k\) in Fourier inversion and multiply by \(q\):
\[
 0=Z+\sum_{r\ne0}\widehat A(r)e_q(rk).
\]
\end{proof}

\begin{remark}[What the identity does and does not say]
\Cref{prop:one-zero} proves that nonzero Fourier modes determine the
zero mode once a structural zero is known.  It does not bound the
coherent sum on the right.  A mean square, a pointwise estimate on
generic \(r\), or an exceptional-set count must still be combined
with the exact phases \(e_q(rk)\).
\end{remark}

The single-zero identity is not optimal when \(A\) has many zero
coordinates.  Let
\[
 S=\supp A,\qquad s=|S|,\qquad K=G\setminus S,\qquad k=|K|=q-s.
\]
Assume \(K\ne\varnothing\).  For a weight \(w:K\to\C\) satisfying
\begin{equation}
 \sum_{x\in K}w(x)=1,
\label{eq:weight-normalization}
\end{equation}
define
\begin{equation}
 m_w(r)=\sum_{x\in K}w(x)e_q(rx).
\label{eq:multiplier}
\end{equation}

\begin{theorem}[Weighted zero-set reconstruction]
\label{thm:weighted-reconstruction}
For every \(w\) satisfying \eqref{eq:weight-normalization},
\begin{equation}
 \boxed{
 Z=-\sum_{r\ne0}m_w(r)\widehat A(r).}
\label{eq:weighted-reconstruction}
\end{equation}
Moreover,
\begin{equation}
 \sum_{r\ne0}|m_w(r)|^2
 =q\sum_{x\in K}|w(x)|^2-1.
\label{eq:multiplier-parseval}
\end{equation}
\end{theorem}

\begin{proof}
Multiply the inversion identity \(A(x)=0\), \(x\in K\), by \(w(x)\)
and sum over \(K\).  The \(r=0\) term is \(Z\) because
\(\sum_Kw=1\); the remaining terms give
\eqref{eq:weighted-reconstruction}.  Parseval applied to the
zero-extension of \(w\) gives
\[
 \sum_{r\in G}|m_w(r)|^2=q\sum_{x\in K}|w(x)|^2.
\]
Since \(m_w(0)=1\), removing \(r=0\) proves
\eqref{eq:multiplier-parseval}.
\end{proof}
